% SEGMENT OLP-0563-S01
\documentclass[../../../include/open-logic-section]{subfiles}

\begin{document}

\olfileid{sth}{spine}{zf}
\olsection{$\Z$ மற்றும் $\ZF$: ஒரு முக்கியக் கட்டம்}

அடித்தளத்துடன் இன்னொரு முக்கியக் கட்டத்தை அடைகிறோம். ஒரு குறிப்பிட்ட
ஒன்று ``நீக்கப்பட்ட'' குறிப்பிட்ட கோட்பாடுகள் என்று நாம் கூறிய
$\Zminus$, $\ZFminus$ ஆகிய கோட்பாடுகளைக் கருதினோம். அந்தக் குறிப்பிட்ட
ஒன்று அடித்தளம். எனவே:

\begin{defn}
$\Zminus$ உடன் அடித்தளத்தை $\Z$ சேர்க்கிறது. ஆகவே அதன் அடிகோள்கள்
உறுப்புசார் சமத்துவம், சேர்ப்பு, ஜோடிகள், அடுக்குக்கணங்கள், முடிவிலி,
அடித்தளம் மற்றும் பிரிப்புத் திட்டத்தின் எல்லா நேர்வுகளும் ஆகும்.

$\ZFminus$ உடன் அடித்தளத்தை $\ZF$ சேர்க்கிறது. வேறுவிதமாகக் கூறினால்,
$\Z$ உடன் மாற்றீட்டின் எல்லா நேர்வுகளையும் $\ZF$ சேர்க்கிறது.
\end{defn}

எனினும், உங்களுக்கு ஒரு கேள்வி தோன்றியிருக்கலாம். $\ZFminus$ தொடர்பில்
ஒழுங்குமை அடித்தளத்திற்குச் சமானமானது என்றால், மேலும் ஒழுங்குமையின்
நியாயப்படுத்தல் தெளிவானது என்றால், ஒழுங்குமையை விட அடித்தளத்தை நமது
அடிப்படை அடிகோளாக ஏற்றுக்கொள்ள ஏன் சுற்றுவழியில் செல்ல வேண்டும்?

வரலாற்றுக் காரணங்களை (யார் எதை எப்போது வகுத்தார் என்பவை) ஒதுக்கினால்,
அடிப்படைக் காரணம் இதுதான்: $V_\alpha$-களின் வரையறையைப் பயன்படுத்தாமலேயே
அடித்தளத்தை முன்வைக்கலாம். அந்த வரையறை \olref[recursion]{sec} இன்
எல்லாப் பணிகளையும் சார்ந்திருந்தது: அது நியாயப்படுத்தப்பட்டது என்று
காட்ட முடிவிலிக்கடந்த மீள்வரையறையை நிறுவ வேண்டியிருந்தது. ஆனால்
முடிவிலிக்கடந்த மீள்வரையறைக்கான நமது நிறுவல் \emph{மாற்றீட்டை}
பயன்படுத்தியது. எனவே அடித்தளமும் ஒழுங்குமையும் $\ZFminus$ தொடர்பில்
சமானமானவை என்றாலும், $\Zminus$ தொடர்பில் அவை சமானமானவை அல்ல.

உண்மையில், இவ்வெளிய குறிப்பைவிட நிலைமை இன்னும் தீவிரமானது. இது இந்நூலின்
எல்லையை வெகுவாக மீறினாலும், $\Zminus$, $\Z$ ஆகிய இரண்டுமே
$V_\alpha$-களை வரையறுக்க முடியாத அளவுக்கு வலிமையற்றவை என்று தெரியவருகிறது.
ஆகவே நீங்கள் $\Z$ இல் மட்டும் பணிபுரிந்தால், ஒழுங்குமை (நாம் அதை
வகுத்துள்ள வடிவில்) \emph{பொருள்கொள்ளவே} இல்லை. இதனால்தான் நமது
அதிகாரபூர்வ அடிகோள் ஒழுங்குமையாக இல்லாமல் அடித்தளமாக உள்ளது.

இனிமேல் வேறுவிதமாகக் கூறப்படாவிட்டால், மேலதிகக் குறிப்பின்றி $\ZF$ இல்
பணிபுரிவோம்.

\end{document}
